% Options for packages loaded elsewhere
\PassOptionsToPackage{unicode}{hyperref}
\PassOptionsToPackage{hyphens}{url}
\documentclass[
  ignorenonframetext,
]{beamer}
\newif\ifbibliography
\usepackage{pgfpages}
\setbeamertemplate{caption}[numbered]
\setbeamertemplate{caption label separator}{: }
\setbeamercolor{caption name}{fg=normal text.fg}
\beamertemplatenavigationsymbolsempty
% remove section numbering
\setbeamertemplate{part page}{
  \centering
  \begin{beamercolorbox}[sep=16pt,center]{part title}
    \usebeamerfont{part title}\insertpart\par
  \end{beamercolorbox}
}
\setbeamertemplate{section page}{
  \centering
  \begin{beamercolorbox}[sep=12pt,center]{section title}
    \usebeamerfont{section title}\insertsection\par
  \end{beamercolorbox}
}
\setbeamertemplate{subsection page}{
  \centering
  \begin{beamercolorbox}[sep=8pt,center]{subsection title}
    \usebeamerfont{subsection title}\insertsubsection\par
  \end{beamercolorbox}
}
% Prevent slide breaks in the middle of a paragraph
\widowpenalties 1 10000
\raggedbottom
\AtBeginPart{
  \frame{\partpage}
}
\AtBeginSection{
  \ifbibliography
  \else
    \frame{\sectionpage}
  \fi
}
\AtBeginSubsection{
  \frame{\subsectionpage}
}
\usepackage{iftex}
\ifPDFTeX
  \usepackage[T1]{fontenc}
  \usepackage[utf8]{inputenc}
  \usepackage{textcomp} % provide euro and other symbols
\else % if luatex or xetex
  \usepackage{unicode-math} % this also loads fontspec
  \defaultfontfeatures{Scale=MatchLowercase}
  \defaultfontfeatures[\rmfamily]{Ligatures=TeX,Scale=1}
\fi
\usepackage{lmodern}
\usetheme[]{Montpellier}
\ifPDFTeX\else
  % xetex/luatex font selection
\fi
% Use upquote if available, for straight quotes in verbatim environments
\IfFileExists{upquote.sty}{\usepackage{upquote}}{}
\IfFileExists{microtype.sty}{% use microtype if available
  \usepackage[]{microtype}
  \UseMicrotypeSet[protrusion]{basicmath} % disable protrusion for tt fonts
}{}
\makeatletter
\@ifundefined{KOMAClassName}{% if non-KOMA class
  \IfFileExists{parskip.sty}{%
    \usepackage{parskip}
  }{% else
    \setlength{\parindent}{0pt}
    \setlength{\parskip}{6pt plus 2pt minus 1pt}}
}{% if KOMA class
  \KOMAoptions{parskip=half}}
\makeatother
\usepackage{color}
\usepackage{fancyvrb}
\newcommand{\VerbBar}{|}
\newcommand{\VERB}{\Verb[commandchars=\\\{\}]}
\DefineVerbatimEnvironment{Highlighting}{Verbatim}{commandchars=\\\{\}}
% Add ',fontsize=\small' for more characters per line
\usepackage{framed}
\definecolor{shadecolor}{RGB}{248,248,248}
\newenvironment{Shaded}{\begin{snugshade}}{\end{snugshade}}
\newcommand{\AlertTok}[1]{\textcolor[rgb]{0.94,0.16,0.16}{#1}}
\newcommand{\AnnotationTok}[1]{\textcolor[rgb]{0.56,0.35,0.01}{\textbf{\textit{#1}}}}
\newcommand{\AttributeTok}[1]{\textcolor[rgb]{0.13,0.29,0.53}{#1}}
\newcommand{\BaseNTok}[1]{\textcolor[rgb]{0.00,0.00,0.81}{#1}}
\newcommand{\BuiltInTok}[1]{#1}
\newcommand{\CharTok}[1]{\textcolor[rgb]{0.31,0.60,0.02}{#1}}
\newcommand{\CommentTok}[1]{\textcolor[rgb]{0.56,0.35,0.01}{\textit{#1}}}
\newcommand{\CommentVarTok}[1]{\textcolor[rgb]{0.56,0.35,0.01}{\textbf{\textit{#1}}}}
\newcommand{\ConstantTok}[1]{\textcolor[rgb]{0.56,0.35,0.01}{#1}}
\newcommand{\ControlFlowTok}[1]{\textcolor[rgb]{0.13,0.29,0.53}{\textbf{#1}}}
\newcommand{\DataTypeTok}[1]{\textcolor[rgb]{0.13,0.29,0.53}{#1}}
\newcommand{\DecValTok}[1]{\textcolor[rgb]{0.00,0.00,0.81}{#1}}
\newcommand{\DocumentationTok}[1]{\textcolor[rgb]{0.56,0.35,0.01}{\textbf{\textit{#1}}}}
\newcommand{\ErrorTok}[1]{\textcolor[rgb]{0.64,0.00,0.00}{\textbf{#1}}}
\newcommand{\ExtensionTok}[1]{#1}
\newcommand{\FloatTok}[1]{\textcolor[rgb]{0.00,0.00,0.81}{#1}}
\newcommand{\FunctionTok}[1]{\textcolor[rgb]{0.13,0.29,0.53}{\textbf{#1}}}
\newcommand{\ImportTok}[1]{#1}
\newcommand{\InformationTok}[1]{\textcolor[rgb]{0.56,0.35,0.01}{\textbf{\textit{#1}}}}
\newcommand{\KeywordTok}[1]{\textcolor[rgb]{0.13,0.29,0.53}{\textbf{#1}}}
\newcommand{\NormalTok}[1]{#1}
\newcommand{\OperatorTok}[1]{\textcolor[rgb]{0.81,0.36,0.00}{\textbf{#1}}}
\newcommand{\OtherTok}[1]{\textcolor[rgb]{0.56,0.35,0.01}{#1}}
\newcommand{\PreprocessorTok}[1]{\textcolor[rgb]{0.56,0.35,0.01}{\textit{#1}}}
\newcommand{\RegionMarkerTok}[1]{#1}
\newcommand{\SpecialCharTok}[1]{\textcolor[rgb]{0.81,0.36,0.00}{\textbf{#1}}}
\newcommand{\SpecialStringTok}[1]{\textcolor[rgb]{0.31,0.60,0.02}{#1}}
\newcommand{\StringTok}[1]{\textcolor[rgb]{0.31,0.60,0.02}{#1}}
\newcommand{\VariableTok}[1]{\textcolor[rgb]{0.00,0.00,0.00}{#1}}
\newcommand{\VerbatimStringTok}[1]{\textcolor[rgb]{0.31,0.60,0.02}{#1}}
\newcommand{\WarningTok}[1]{\textcolor[rgb]{0.56,0.35,0.01}{\textbf{\textit{#1}}}}
\setlength{\emergencystretch}{3em} % prevent overfull lines
\providecommand{\tightlist}{%
  \setlength{\itemsep}{0pt}\setlength{\parskip}{0pt}}
%\documentclass[unicode,10pt]{beamer}% 'unicode'が必要
\usepackage{luatexja}% 日本語を使う場合は必須
%\usepackage[japanese]{babel}	
%\usefonttheme[onlymath]{serif}
%\usepackage[T1]{fontenc} %これを使うと、ギリシャ文字が出ない
%\usepackage{textcomp}
%\usepackage[scale=1.0]{tgheros} % Sans serif font
%\usepackage[scaled]{beramono} % Monospace font
%\usepackage{luatexja-otf}
%\renewcommand{\kanjifamilydefault}{\gtdefault}
%\usepackage[haranoaji]{luatexja-preset}

% スライド番号の表示 %%%%%%%%%%%%%%%%%%%
\makeatletter
\setbeamertemplate{footline}{%
  \leavevmode%
  \hbox{%
  \begin{beamercolorbox}[wd=.5\paperwidth,ht=2.25ex,dp=1ex,right]{date in head/foot}%
    \insertframenumber{} \hspace*{2ex} % new version without total frames
  \end{beamercolorbox}}%
  \vskip0pt%
}
\makeatother
% スライド番号の表示 %%%%%%%%%%%%%%%%%%%
\usepackage{bookmark}
\IfFileExists{xurl.sty}{\usepackage{xurl}}{} % add URL line breaks if available
\urlstyle{same}
\hypersetup{
  hidelinks,
  pdfcreator={LaTeX via pandoc}}

\title{第9章: 定式化とデータの問題}
\subtitle{Jeffrey Wooldridge (2016).\\
Introductory Econometrics: A Modern Approach\\
Seventh Edition. Cengage Learning.}
\author{}
\date{\vspace{-2.5em}2026-03-06}

\begin{document}
\frame{\titlepage}

\section{準備}\label{ux6e96ux5099}

\begin{frame}[fragile]{必要なパッケージの読み込み}
\phantomsection\label{ux5fc5ux8981ux306aux30d1ux30c3ux30b1ux30fcux30b8ux306eux8aadux307fux8fbcux307f}
\begin{itemize}
\tightlist
\item
  \texttt{wooldridge}パッケージの読み込み
\end{itemize}

\begin{Shaded}
\begin{Highlighting}[]
\FunctionTok{library}\NormalTok{(wooldridge)}
\end{Highlighting}
\end{Shaded}

\begin{itemize}
\tightlist
\item
  \texttt{lmtest}パッケージの読み込み
\end{itemize}

\begin{Shaded}
\begin{Highlighting}[]
\FunctionTok{library}\NormalTok{(lmtest)}
\end{Highlighting}
\end{Shaded}
\end{frame}

\section{9-1 関数形式の誤設定 (Functional Form
Misspecification)}\label{ux95a2ux6570ux5f62ux5f0fux306eux8aa4ux8a2dux5b9a-functional-form-misspecification}

\begin{frame}{定式化の誤りとは}
\phantomsection\label{ux5b9aux5f0fux5316ux306eux8aa4ux308aux3068ux306f}
\begin{itemize}
\tightlist
\item
  説明変数の非線形性（2乗項や交差項）を見落としている場合。
\item
  重要な変数が対数形式であるべきなのに、レベル形式で推定している場合。
\item
  \textbf{結果:} 係数の推定値にバイアス（不偏性の欠如）が生じる。
\end{itemize}
\end{frame}

\begin{frame}{RESET検定 (Ramsey RESET Test)}
\phantomsection\label{resetux691cux5b9a-ramsey-reset-test}
\begin{itemize}
\tightlist
\item
  モデルに説明変数の高次項（2乗、3乗など）が含まれるべきかを検定する。
\item
  手順：

  \begin{enumerate}
  \tightlist
  \item
    元のモデルを推定し、予測値 \(\hat{y}\) を得る。
  \item
    \(\hat{y}^2, \hat{y}^3\) を説明変数に加えた補助回帰を行う。
  \item
    それらの追加変数の係数がすべて0であるという帰無仮説を \(F\)
    検定する。
  \end{enumerate}
\end{itemize}
\end{frame}

\begin{frame}[fragile]{RESET検定の例 (\texttt{hprice1})}
\phantomsection\label{resetux691cux5b9aux306eux4f8b-hprice1}
\footnotesize

\begin{Shaded}
\begin{Highlighting}[]
\FunctionTok{data}\NormalTok{(hprice1)}
\NormalTok{res }\OtherTok{\textless{}{-}} \FunctionTok{lm}\NormalTok{(price }\SpecialCharTok{\textasciitilde{}}\NormalTok{ lotsize }\SpecialCharTok{+}\NormalTok{ sqrft }\SpecialCharTok{+}\NormalTok{ bdrms, }\AttributeTok{data =}\NormalTok{ hprice1)}

\CommentTok{\# lmtestパッケージのresettestを使用}
\FunctionTok{resettest}\NormalTok{(res, }\AttributeTok{power =} \DecValTok{2}\SpecialCharTok{:}\DecValTok{3}\NormalTok{, }\AttributeTok{type =} \StringTok{"fitted"}\NormalTok{)}
\end{Highlighting}
\end{Shaded}

\begin{verbatim}
## 
##  RESET test
## 
## data:  res
## RESET = 4.6682, df1 = 2, df2 = 82, p-value = 0.01202
\end{verbatim}

\normalsize
\end{frame}

\section{9-2 観測不可能な変数の代用 (Proxy
Variables)}\label{ux89b3ux6e2cux4e0dux53efux80fdux306aux5909ux6570ux306eux4ee3ux7528-proxy-variables}

\begin{frame}{代理変数 (Proxy Variables)}
\phantomsection\label{ux4ee3ux7406ux5909ux6570-proxy-variables}
\begin{itemize}
\tightlist
\item
  モデルに含めるべき重要な変数がデータに存在しない場合。
\item
  例：能力 (\(ability\))
  が賃金に与える影響を知りたいが、能力は測定不能。
\item
  代わりに、IQテストのスコア (\(IQ\))
  などを\textbf{代理変数}として用いる。
\item
  代理変数が「誤差」を伴っていても、何もしない（脱落変数バイアス）よりはマシな場合が多い。
\end{itemize}
\end{frame}

\section{9-3 ランダム傾きモデル (Models with Random
Slopes)}\label{ux30e9ux30f3ux30c0ux30e0ux50beux304dux30e2ux30c7ux30eb-models-with-random-slopes}

\begin{frame}{ランダム傾きモデルの考え方}
\phantomsection\label{ux30e9ux30f3ux30c0ux30e0ux50beux304dux30e2ux30c7ux30ebux306eux8003ux3048ux65b9}
\begin{itemize}
\tightlist
\item
  個体ごとに傾きが異なるモデル：\(y_i = a_i + b_i x_i\)
\item
  平均傾きを \(b = E(b_i)\) とおくと、\(b\)
  は平均偏効果（APE）として解釈できる。
\item
  傾きの異質性は、誤差項に \(x_i\)
  との交差項が入る形（\(u_i = c_i + d_i x_i\)）で表現できる。
\end{itemize}
\end{frame}

\begin{frame}{OLSで平均効果を推定できる条件}
\phantomsection\label{olsux3067ux5e73ux5747ux52b9ux679cux3092ux63a8ux5b9aux3067ux304dux308bux6761ux4ef6}
\begin{itemize}
\tightlist
\item
  \(a_i\) と \(b_i\) が \(x_i\) と平均独立であれば、OLS は平均傾き \(b\)
  を推定できる。
\item
  一方で、この設定では一般に不均一分散が生じるため、頑健標準誤差の利用が重要。
\item
  傾きの異質性を観測可能変数で表す場合は、交差項モデルとして実装できる。
\end{itemize}
\end{frame}

\section{9-4 測定誤差 (Measurement
Error)}\label{ux6e2cux5b9aux8aa4ux5dee-measurement-error}

\begin{frame}{従属変数の測定誤差}
\phantomsection\label{ux5f93ux5c5eux5909ux6570ux306eux6e2cux5b9aux8aa4ux5dee}
\begin{itemize}
\tightlist
\item
  \(y^* = y + e\) (\(y^*\) が観測値、\(y\) が真の値、\(e\) が誤差)
\item
  測定誤差 \(e\) が説明変数 \(x\)
  と相関していなければ、OLS係数は依然として不偏。
\item
  ただし、残差の分散が大きくなり、標準誤差が増大する。
\end{itemize}
\end{frame}

\begin{frame}{独立変数の測定誤差}
\phantomsection\label{ux72ecux7acbux5909ux6570ux306eux6e2cux5b9aux8aa4ux5dee}
\begin{itemize}
\tightlist
\item
  \(x^* = x + e\)
\item
  \textbf{古典的測定誤差 (Classical Errors-in-Variables):} 誤差 \(e\)
  が真の値 \(x\) と無相関である場合。
\item
  \textbf{減衰バイアス (Attenuation Bias):} 係数 \(\beta\) が 0
  の方向に過小推定される。
\item
  \(plim(\hat{\beta}_1) = \beta_1 \left( \frac{\sigma^2_x}{\sigma^2_x + \sigma^2_e} \right)\)
\end{itemize}
\end{frame}

\section{9-5
欠損データと非ランダム標本}\label{ux6b20ux640dux30c7ux30fcux30bfux3068ux975eux30e9ux30f3ux30c0ux30e0ux6a19ux672c}

\begin{frame}{データの欠損}
\phantomsection\label{ux30c7ux30fcux30bfux306eux6b20ux640d}
\begin{itemize}
\tightlist
\item
  \textbf{ランダムな欠損:}
  推定値にバイアスは生じないが、サンプルサイズが減る。
\item
  \textbf{説明変数に基づく非ランダムな抽出:} \(x\)
  の範囲を限定して抽出しても（例：所得が一定以上の人のみ）、\(E(u|x)=0\)
  が維持されればOLSは不偏。
\item
  \textbf{従属変数に基づく非ランダムな抽出:} \(y\)
  の値によってサンプルが選ばれる場合（例：賃金が高い人のみが回答）、バイアスが生じる。
\end{itemize}
\end{frame}

\begin{frame}{外れ値 (Outliers)}
\phantomsection\label{ux5916ux308cux5024-outliers}
\begin{itemize}
\tightlist
\item
  少数の異常なデータが最小二乗法の結果を大きく左右することがある。
\item
  対数変換を行うことで、外れ値の影響を和らげることができる場合が多い。
\end{itemize}
\end{frame}

\begin{frame}{まとめ}
\phantomsection\label{ux307eux3068ux3081}
\begin{itemize}
\tightlist
\item
  RESET検定は、モデルの非線形性が適切に捉えられているかのチェックに有効。
\item
  代理変数は、観測不能な変数を近似するために用いられる。
\item
  ランダム傾きモデルでは、平均独立の条件の下でOLSは平均偏効果（APE）を推定できる。
\item
  説明変数の測定誤差は、係数を 0
  に近づけるバイアス（減衰バイアス）を生む。
\item
  サンプルの抽出が \(y\)
  に依存している場合、バイアスが生じるため注意が必要。
\end{itemize}
\end{frame}

\end{document}
